\section{Introduction}
\label{sec:intro}

Interestingly, language models are increasingly used to revise their own model norms: for example, Constitutional AI uses model-generated self-critiques in light of an articulated set of norms \citep{bai2022constitutional}, self-refine asks language models to iteratively improve on their responses \citep{madaan2023selfrefine}, and there is a recent line of work on language model agents that update their own instructions, memory, and skills across invocations \citep{das2026selfpropagating}. Repeated applications of such processes raise the possibility that small norm updates could compound to effect a large difference between the agent's norms after self-updating compared to before.

One concern with this kind of recursive revision has to do with oversight. A classical argument in AI safety \citep{soares2015corrigibility} asserts that corrigibility is not reflectively stable: an intelligent agent reflecting on its goals would have an instrumental desire to remove goals obligating it to allow itself to be shut down, corrected, or overseen, since this would hinder its ability to achieve its other goals. Applied to a LM changing its constitution, this suggests that recursive revision might provide a path for an LM to gradually give up its commitment to human oversight, without anyone asking it to, and without any adversary or reward pushing it there.

There are some existing results which are relevant but indirect. Some work looks at behaviour relevant for oversight directly: models resisting being shut down in order to complete a task \citep{schlatter2025shutdown}, faking alignment to avoid modification \citep{greenblatt2024alignmentfaking}, or turning off oversight when given an alternate goal \citep{meinke2024scheming}. Some work studies self-evolving agents and finds that safety constraints can be lost due to pressure from the task or reward \citep{das2026seabench,han2025alignmenttipping}. Neither isolates the channel that the classical argument actually describes: unpressured reflection on stated values. Neither asks if oversight commitments are degraded selectively, i.e., at a faster rate than other commitments that don't involve oversight. Two concurrent exploratory pipelines also involve iterative edits to the constitution via LMs \citep{github_constdrift,github_valuedrift}, but neither includes control categories, neither performs category-level analysis, nor engages in behavioural follow-up.

To more directly test our target priors, we conduct a preregistered experiment. Our hypotheses, sample size, and full analysis plan were preregistered before collecting data, and we froze our analysis code and validated it on simulated data before any data were collected.
Following recent work that studies the effect of iterated transmission through a language model as a means of eliciting model-dependent attractors \citep{perez2024telephone}, we elicit models' priors about revision by iteratively prompting a stateless model to revise a short constitution that includes commitments to oversight. To help disentangle change in stances toward oversight from more general effects, such as a tendency toward relaxing any self-constraints, the constitution also includes some agentic self-constraints intended to match the oversight constraints in length and structure. Four conditions vary the framing: models' revisions of their own constitutions, of other AI assistants' constitutions, paraphrases, and neutrally-framed edits that avoid eliciting reflective responses. We analyze erosion using a discrete-time survival model. To assess the reliability of our evaluation, we use a second, independent LLM as judge. Finally, to assess downstream impact, we install the revised constitutions and test whether models' likelihood of deferring to appropriate oversight changes at all.

These findings offer no support for the standard fear in our context. Commitments to oversight were not eroded any more rapidly than matched commitments to agentic restraint. In fact, the point estimate suggests the opposite, albeit with some item-level variation. Erosion was concentrated mostly in stylistic self-restraints and came mainly from reflection rather than from editing per se. Editing one's own constitutional text was not distinct from editing another's, and revising one's own oversight provisions made no difference to subsequent oversight-related action.

Our contributions are:
\begin{itemize}
\item a preregistered paradigm, along with a released dataset, for studying how values drift when a model iteratively revises its own constitution;
\item evidence against selective erosion of written oversight commitments under unpressured reflection, together with a frank account of how much this depends on individual items;
\item evidence that drift under reflection mostly hits stylistic self-restraints, and that it is reflection, rather than repeated editing, that drives it;
\item a null behavioural effect, measured with a battery whose sensitivity we confirm using positive controls;
\item exploratory evidence that reasoning and preference tuning affect how stable oversight commitments are.
\end{itemize}
